\providecommand{\Res}[2]{\qty[round-mode=figures,round-precision=3]{\fpeval{#1}}{#2}}
\providecommand{\Val}[1]{\num[round-mode=figures,round-precision=4]{\fpeval{#1}}}
\subsubsection*{Question a) Normale au sommet de la bosse}
\begin{enumerate}[label=\arabic*)]
    \item \textbf{Variables connues et recherchées :}
    \begin{itemize}[label= -]
        \item Masse de la voiture : $m = \qty{\v{m}}{kg}$
        \item Vitesse de la voiture : $v = \qty{\v{vk}}{km/h}$ $= \qty[round-mode=figures,round-precision=4,drop-zero-decimal]{\vsi{vk}}{m/s}$
        \item Rayon de courbure : $r = \qty{\v{r}}{m}$
        \item Le centre du cercle est sous la voiture : $mg$ vers le centre, $n$ vers l'extérieur
        \item Variable recherchée : $n$ (N)
    \end{itemize}
    \item \textbf{Équations pertinentes :}
    \[ mg - n = \frac{m v^2}{r} \]
    \item \textbf{Isolation de la variable recherchée :}
    \[ n = m\left(g - \frac{v^2}{r}\right) \]
    \item \textbf{Résolution numérique :}
    \[ n = \num{\vsi{m}}\left(\num{9.81} - \frac{(\Val{\vsi{vk}})^2}{\num{\vsi{r}}}\right) = \Res{\vsi{m}*(9.81 - \vsi{vk}^2/\vsi{r})}{N} \]
    \item \textbf{Réponse :} la normale vaut \textbf{\Res{\vsi{m}*(9.81 - \vsi{vk}^2/\vsi{r})}{N}}, moins que le poids : on se sent plus léger.
\end{enumerate}

\subsubsection*{Question b) Normale au fond du creux}
\begin{enumerate}[label=\arabic*)]
    \item \textbf{Variables connues et recherchées :}
    \begin{itemize}[label= -]
        \item Masse de la voiture : $m = \qty{\v{m}}{kg}$
        \item Vitesse de la voiture : $v = \qty{\v{vk}}{km/h}$ $= \qty[round-mode=figures,round-precision=4,drop-zero-decimal]{\vsi{vk}}{m/s}$
        \item Rayon de courbure : $r = \qty{\v{r}}{m}$
        \item Le centre est au-dessus : $n$ vers le centre, $mg$ vers l'extérieur
        \item Variable recherchée : $n$ (N)
    \end{itemize}
    \item \textbf{Équations pertinentes :}
    \[ n - mg = \frac{m v^2}{r} \]
    \item \textbf{Isolation de la variable recherchée :}
    \[ n = m\left(g + \frac{v^2}{r}\right) \]
    \item \textbf{Résolution numérique :}
    \[ n = \num{\vsi{m}}\left(\num{9.81} + \frac{(\Val{\vsi{vk}})^2}{\num{\vsi{r}}}\right) = \Res{\vsi{m}*(9.81 + \vsi{vk}^2/\vsi{r})}{N} \]
    \item \textbf{Réponse :} la normale vaut \textbf{\Res{\vsi{m}*(9.81 + \vsi{vk}^2/\vsi{r})}{N}} : on se sent plus lourd.
\end{enumerate}

\subsubsection*{Question c) Vitesse de décollage}
\begin{enumerate}[label=\arabic*)]
    \item \textbf{Variables connues et recherchées :}
    \begin{itemize}[label= -]
        \item Rayon de courbure : $r = \qty{\v{r}}{m}$
        \item Contact perdu quand $n = 0$
        \item Variable recherchée : $v$ (m/s)
    \end{itemize}
    \item \textbf{Équations pertinentes :}
    \[ mg = \frac{m v^2}{r} \]
    \item \textbf{Isolation de la variable recherchée :}
    \[ v = \sqrt{g r} \]
    \item \textbf{Résolution numérique :}
    \[ v = \sqrt{\num{9.81}\cdot\num{\vsi{r}}} = \Res{sqrt(9.81*\vsi{r})}{m/s} = \Res{sqrt(9.81*\vsi{r})*3.6}{km/h} \]
    \item \textbf{Réponse :} au-delà de \textbf{\Res{sqrt(9.81*\vsi{r})}{m/s}}, la voiture quitte la route au sommet.
\end{enumerate}
